\documentclass[a4paper,12pt]{article}

\usepackage[utf8]{inputenc}
\usepackage[T1]{fontenc}
\usepackage{amssymb}
\usepackage{amsmath}
\usepackage{epigraph}						% Funny quotes
\usepackage{bm}								% Bold face for vectors/matrices
\usepackage[hidelinks]{hyperref}			% Invisible, clickable links in pdf
\usepackage{listings}						% Display source code nicely
\usepackage{mcode}							% Ditto MATLAB code
\usepackage{algpseudocode}					% Nice Al Gore Rhythms
\usepackage{graphicx}						% Including graphics
\usepackage{epstopdf}						% Self-explanatory
\usepackage{float} 							% Better float placing (H option)
\usepackage{siunitx}						% SI units
\usepackage{subcaption}						% Puts stuff next to each otter
\usepackage{newclude}						% \include without clearpage
\bibliographystyle{ieeetr}					% Citations in IEEE style
\lstset{frame=topbottom}

\newcommand{\M}[1]{\bm{#1}} 				% Use for variable matrices
\newcommand{\Mc}[1]{\mathbf{#1}} 			% Use for constant matrices (I, etc.)
\newcommand{\V}[1]{\mathbf{#1}} 			% Use for vectors
\newcommand{\transpose}{^{\text{T}}} 		% Make a guess, bitch
\newcommand{\sub}[1]{_{\textnormal{#1}}} 	% For upright subscripts

\lstset{
    frame=single,
    breaklines=true,
    postbreak=\raisebox{0ex}[0ex][0ex]{\ensuremath{\color{green}\hookrightarrow\space}},
    morekeywords={model, import, constant, parameter, equation, end},
    keywordstyle=\color{blue},
    numbers=left,
    numberstyle=\scriptsize,
    stepnumber=1,
    numbersep=8pt
}

\begin{document}

\begin{titlepage}
\begin{center}

\textsc{\LARGE TTK4135}\\[1.5cm]
\textsc{\LARGE Lab Report} \linebreak
{ \huge \bfseries \textcolor[rgb]{0,0,1}{Real-world application of optimal control methods}\\[0.5cm] }
\today \vspace{1cm} \linebreak 
\textbf{Group 9} \linebreak \linebreak
\emph{\textbf{Participants:}} \linebreak 
712623, 742401, 742434, 742441, 742446 \linebreak


\begin{abstract}

This is a report for the mandatory lab exercise in the course TTK4135 Optimization and Control held at NTNU the spring 2015. The purpose of this exercise is to get practical experience in formulating dynamic optimization problems, discretise them and solve them using a computer, as well as using these results to implement optimal controllers with and without feedback. The optimization problems that got solved were quadratic minimization problems, and the feedback was introduced using LQ controllers.

\end{abstract}

\vfill

\textsc{Faculty of Information Technology, Mathematics and Electrical Engineering}\linebreak
\textsc{Department of Engineering Cybernetics} 

\vspace{1cm}
\includegraphics[width=200pt, keepaspectratio=true]{logo_ntnu_eng.eps}

\end{center}
\end{titlepage}

\tableofcontents

\section{Introduction}

This is a lab called the Helicopter Lab. As the name indicates we are going to control a model of a helicopter in different ways. The system with its hardware is described in detail in \cite{_helicopter_2015}. A model of the helicopter with regulators is allready made for us. The objective is for us, the students, to find an optimal sequence of inputs that will steer the helicopter on the desired path. This will be done by formulating optimization problems, and solving them. We will solve them as discrete problems using MATLAB, wich mean that we first have to discretize our system. This input sequense we give to our system using Simulink/QuaRC. We will run our helicopters using the optimal input directly, and them compare the behaviour when we introduce feedback.

In this report we will be presenting our results in the same order that we obtained them, meaning we will have one section for each exercise in \cite{_helicopter_2015}, where the relevant math, code, plots, results, etc. will be added. The discussion will be done consecutively, with the exeption of one summarizing discution at the end of the report. 
\include*{10.1}
\include*{10.2}
\include*{10.3}
\include*{10.4}
\include*{discussionAndConclusion}

\bibliography{ref-helikopter}
\section{Appendix}
\subsection{Variables}
\begin{center}
\begin{tabular}{l l l l}
\hline 
Symbol & Variable \\ \hline
$p$ & Pitch \\
$p_c$ & Pitch setpoint \\
$\lambda$ & Travel \\
$r$ & Travel rate \\
$r_c$ & Travel rate setpoint \\
$e$ & Elevation \\
$e_c$ & Elevation setpoint \\
$V_f$ & Voltage, front motor \\
$V_b$ & Voltage, rear motor \\
$V_d$ & Voltage difference, $V_f -V_b$ \\
$V_s$ & Voltage sum, $V_f + V_b$ \\
$K_{pp}$, $K_{pd}$, $K_{ep}$, $K_{ei}$, $K_{ed}$ & Controller gains \\
$T_g$ & Moment required to keep the helicopter flying
\end{tabular}
\end{center}
\label{table:variables}



\vspace{1cm}

\subsection{Constants}
\begin{center}
\begin{tabular}{l l l l}
\hline
Symbol & Parameter & Value & Unit \\ \hline
$l_a$ & Distance from elevation axis to helicopter body & 0.64 & m \\
$l_h$ & Distance from pitch axis to motor & 0.177 & m \\
$K_f$ & Distance from elevation axis to helicopter body & 0.1983 &\si{\newton \per \meter}  \\
$J_e$ & Moment of inertia (elevation axis) & 1.0625 & \si{\kilogram \meter \squared}  \\
$J_t$ & Moment of inertia (travel axis) & 1.0625 & \si{\kilogram \meter \squared}  \\
$J_p$ & Moment of inertia (pitch axis) & 0.0406 & \si{\kilogram \meter \squared}  \\
$m_h$ & Mass of helicopter & 1.297 &\si{\kilogram}  \\
$m_w$ & Counterweight & 1.802 & \si{\kilogram} \\
$m_g$ & Effective mass of the helicopter & 0.026 & \si{\kilogram} \\
$K_p$ & Force to lift the helicopter from the ground & 0.2551 & \si{\newton}
\end{tabular}
\end{center}
\label{table:constants}

\subsection{Code for (10.2) and (10.3)}

This code is modified from the code handed out on itslearning.

\begin{lstlisting}
init04;
delta_t	  = 0.25;	                    % sampling time
sek_forst = 5;

q = 0.1;
% System model. x=[lambda r p p_dot]'

A1 = [1 delta_t 0 0;
      0 1 -K_2*delta_t 0;
      0 0 1 delta_t;
      0 0 -K_1*K_pp*delta_t 1-K_1*K_pd*delta_t];
  
B1 = [0; 0; 0; K_1*K_pp*delta_t];

% Number of states and inputs

mx = size(A1,2);                        % Number of states (number of columns in A)
mu = size(B1,2);                        % Number of inputs(number of columns in B)

% Initial values

x1_0 = pi;                              % Lambda
x2_0 = 0;                               % r
x3_0 = 0;                               % p
x4_0 = 0;                               % p_dot
x0   = [x1_0 x2_0 x3_0 x4_0]';          % Initial values

% Time horizon and initialization

N  = 100;                                % Time horizon for states
M  = N;                                 % Time horizon for inputs
z  = zeros(N*mx+M*mu,1);                % Initialize z for the whole horizon
z0 = z;                                 % Initial value for optimization

% Bounds

ul 	    = -30*pi/180;                   % Lower bound on control -- u1
uu 	    = 30*pi/180;                    % Upper bound on control -- u1

xl      = -Inf*ones(mx,1);              % Lower bound on states (no bound)
xu      = Inf*ones(mx,1);               % Upper bound on states (no bound)
xl(3)   = ul;                           % Lower bound on state x3
xu(3)   = uu;                           % Upper bound on state x3

% Generate constraints on measurements and inputs

[vlb,vub]       = genBegr2(N,M,xl,xu,ul,uu);
vlb(N*mx+M*mu)  = 0;                    % We want the last input to be zero
vub(N*mx+M*mu)  = 0;                    % We want the last input to be zero

% Generate the matrix Q and the vector c (objecitve function weights in the QP problem) 

Q1 = zeros(mx,mx);
Q1(1,1) = 1;                            % Weight on state x1
Q1(2,2) = 0;                            % Weight on state x2
Q1(3,3) = q;                            % Weight on state x3
Q1(4,4) = 0;                            % Weight on state x4
P1 = q;                                 % Weight on input
Q = 2*genq2(Q1,P1,N,M,mu);              % Generate Q
c = zeros(N*mx+M*mu,1);                 % Generate c

% Generate system matrixes for linear model

Aeq = gena2(A1,B1,N,mx,mu);          % Generate A
beq = zeros(1, size(Aeq,1));         % Generate b
beq(1:mx) = A1*x0; 	        	        % Initial value

% Solve Qp problem with linear model
tic
[z,lambda] = quadprog(Q, c, [], [], Aeq, beq, vlb, vub, z0);
t1=toc;


% Extract control inputs and states

u  = [z(N*mx+1:N*mx+M*mu);z(N*mx+M*mu)]; % Control input from solution

x1 = [x0(1);z(1:mx:N*mx)];              % State x1 from solution
x2 = [x0(2);z(2:mx:N*mx)];              % State x2 from solution
x3 = [x0(3);z(3:mx:N*mx)];              % State x3 from solution
x4 = [x0(4);z(4:mx:N*mx)];              % State x4 from solution

Antall = 5/delta_t;
Nuller = zeros(Antall,1);
Enere  = ones(Antall,1);

u   = [Nuller; u; Nuller];
x1  = [pi*Enere; x1; Nuller];
x2  = [Nuller; x2; Nuller];
x3  = [Nuller; x3; Nuller];
x4  = [Nuller; x4; Nuller];

%save trajektor1ny
t = 0:delta_t:delta_t*(length(u)-1); % real time

simin = [t' u];

% figure
                

figure(2)
subplot(511)
stairs(t,u),grid
ylabel('u')
subplot(512)
plot(t,x1,'m',t,x1,'mo'),grid
ylabel('lambda')
subplot(513)
plot(t,x2,'m',t,x2','mo'),grid
ylabel('r')
subplot(514)
plot(t,x3,'m',t,x3,'mo'),grid
ylabel('p')
subplot(515)
plot(t,x4,'m',t,x4','mo'),grid
xlabel('tid (s)'),ylabel('pdot')

x = [t' x1 x2 x3 x4];
%%

Q_k = [25 0 0 0;
       0 0.5 0 0;
       0 0 100 0;
       0 0 0 0.5];
R = 1;

[K S E] = dlqr(A1, B1, Q_k, R, 0);

\end{lstlisting}

\subsection{Code for (10.4)}

This code is modified from the code handed out on itslearning.

\begin{lstlisting}
init04;
delta_t	  = 0.25;	                    % sampling time
sek_forst = 5;

q = 0.1;
% System model. x=[lambda r p p_dot]'

A = [1 delta_t 0 0 0 0;
      0 1 -K_2*delta_t 0 0 0;
      0 0 1 delta_t 0 0;
      0 0 -K_1*K_pp*delta_t 1-K_1*K_pd*delta_t 0 0;
      0 0 0 0 1 delta_t;
      0 0 0 0 -delta_t*K_3*K_ep 1-delta_t*K_3*K_ed];
  
B = [0 0; 0 0; 0 0; K_1*K_pp*delta_t 0; 0 0; 0 delta_t*K_3*K_ep];

% Number of states and inputs

mx = size(A,2);                        % Number of states (number of columns in A)
mu = size(B,2);                        % Number of inputs(number of columns in B)

% Initial values

x1_0 = pi;                              % Lambda
x2_0 = 0;                               % r
x3_0 = 0;                               % p
x4_0 = 0;                               % p_dot
x5_0 = 0;                               % e
x6_0 = 0;                               % e_dot
x0   = [x1_0 x2_0 x3_0 x4_0 x5_0 x6_0]';          % Initial values

u1_0 = 0;
u2_0 = 0; % ask studass

q_1 = 1;
q_2 = 1;
alfa = 0.2;
beta = 20;
lambda_t = 2*pi/3;

% Time horizon and initialization

N  = 40;                                % Time horizon for states
M  = N;                                 % Time horizon for inputs
z  = zeros(N*mx+M*mu,1);                % Initialize z for the whole horizon
z0 = z;                                 % Initial value for optimization


% Bounds

ul 	    = -30*pi/180;                   % Lower bound on control -- u1
uu 	    = 30*pi/180;                    % Upper bound on control -- u1

xl      = -Inf*ones(mx,1);              % Lower bound on states (no bound)
xu      = Inf*ones(mx,1);               % Upper bound on states (no bound)
xl(3)   = ul;                           % Lower bound on state x3
xu(3)   = uu;                           % Upper bound on state x3

% Generate constraints on measurements and inputs

[vlb,vub]       = genBegr2(N,M,xl,xu,ul,uu);
vlb(N*mx+M*mu)  = 0;                    % We want the last input to be zero
vub(N*mx+M*mu)  = 0;                    % We want the last input to be zero

% Generate system matrixes for linear model
Aeq = gena2(A,B,N,mx,mu);          % Generate A
beq = zeros(1, size(Aeq,1));         % Generate b
beq(1:mx) = A*x0; 	        	        % Initial value

% Generate the matrix Q and the vector c (objecitve function weights in the QP problem) 

Q1 = zeros(mx,mx);
Q1(1,1) = 1;                            % Weight on state x1
Q1(2,2) = 0;                            % Weight on state x2
Q1(3,3) = q;                            % Weight on state x3
Q1(4,4) = 0;                            % Weight on state x4
Q1(5,5) = 0;
Q1(6,6) = 0;
P1 = zeros(mu, mu);
P1(1,1) = q_1;                            % Weight on input
P1(2,2) = q_2;
Q = 2*genq2(Q1,P1,N,M,mu);              % Generate Q
c = zeros(N*mx+M*mu,1);                 % Generate c

% Solve Qp problem with linear model

costf = @(z) 0.5*z'*Q*z;
%tic
%[z,lambda] = quadprog(Q, c, [], [], Aeq, beq, vlb, vub, z0);
%t1=toc;
nonlcon = @(z) alfa*exp(-beta*(z(1+mx)-lambda_t)^2)-z(5+mx);
z = fmincon(costf, z0, [], [], Aeq, beq, vlb, vub, @mycon);
% Calculate objective value

phi1 = 0.0;
PhiOut = zeros(N*mx+M*mu,1);
for i=1:N*mx+M*mu
  phi1=phi1+Q(i,i)*z(i)*z(i);
  PhiOut(i) = phi1;
end

% Extract control inputs and states

u1 = [u1_0; z(N*mx+1:mu:N*mx+M*mu)]; % Control input from solution
u2 = [u2_0; z(N*mx+2:mu:N*mx+M*mu)];

x1 = [x0(1);z(1:mx:N*mx)];              % State x1 from solution
x2 = [x0(2);z(2:mx:N*mx)];              % State x2 from solution
x3 = [x0(3);z(3:mx:N*mx)];              % State x3 from solution
x4 = [x0(4);z(4:mx:N*mx)];              % State x4 from solution
x5 = [x0(5);z(5:mx:N*mx)];
x6 = [x0(6);z(6:mx:N*mx)];

Antall = 5/delta_t;
Nuller = zeros(Antall,1);
Enere  = ones(Antall,1);

u1   = [Nuller; u1; Nuller];
u2   = [Nuller; u2; Nuller];

x1  = [x1_0*Enere; x1; Nuller];
x2  = [Nuller; x2; Nuller];
x3  = [Nuller; x3; Nuller];
x4  = [Nuller; x4; Nuller];
x5 = [Nuller; x5; Nuller];
x6 = [Nuller; x6; Nuller];

%save trajektor1ny
t = 0:delta_t:delta_t*(length(u1)-1); % real time

pc = [t' u1];
ec = [t' u2];
% figure
                

figure(2)
subplot(8,1,1)
stairs(t,u1),grid
ylabel('u1')
subplot(8,1,2)
stairs(t, u2),grid
ylabel('u2')
subplot(8,1,3)
plot(t,x1,'m',t,x1,'mo'),grid
ylabel('lambda')
subplot(8,1,4)
plot(t,x2,'m',t,x2','mo'),grid
ylabel('r')
subplot(8,1,5)
plot(t,x3,'m',t,x3,'mo'),grid
ylabel('p')
subplot(8,1,6)
plot(t,x4,'m',t,x4','mo'),grid
ylabel('pdot')
subplot(8,1,7)
plot(t,x5,'m',t,x5','mo'),grid
ylabel('e')
subplot(8,1,8)
plot(t,x6,'m',t,x6','mo'),grid
ylabel('edot')
xlabel('tid (s)')

x = [t' x1 x2 x3 x4 x5 x6];

%%
Q_k = [2 0 0 0 0 0;
       0 1 0 0 0 0;
       0 0 1 0 0 0;
       0 0 0 1.2 0 0;
       0 0 0 0 1 0;
       0 0 0 0 0 1.5];

R = [1 0;
     0 1];

[K S E] = dlqr(A, B, Q_k, R, 0);

\end{lstlisting}

\end{document}
